% answers.tex -- answers and worked solutions for booklet.tex (Module 5 MC booklet), same
% layout (quiz.sty). One practiceans per quizq, in the same order (72).
% Every number: checks/m5-mc-booklet/numbers.py
% Build: TEXINPUTS=.: xelatex answers.tex (twice)
\documentclass[12pt]{article}
\usepackage{quiz}
\hypersetup{pdftitle={Multiple Choice Booklet - Module 5 - Advanced Mechanics - Answers}}
% links between the two PDF files are not possible, so question headings are plain text here
\renewcommand\hyperlink[2]{#2}
\newcommand\anstopic[1]{\par\Needspace{8\baselineskip}\vspace{4mm}%
  \begin{center}{\headfont\bfseries\fontsize{14}{18}\selectfont #1\par}\end{center}\vspace{2pt}}

\begin{document}
\quiztitle{Year 12 Physics}{Multiple Choice Booklet - Module 5 - Advanced Mechanics: Answers}

\begin{center}
  \renewcommand\arraystretch{1.45}
  \begin{tabular}{|l|*{12}{>{\centering\arraybackslash}p{7mm}|}}
    \hline
    \textbf{Q} & 1 & 2 & 3 & 4 & 5 & 6 & 7 & 8 & 9 & 10 & 11 & 12\\ \hline
    \textbf{Ans} & B & A & C & A & C & B & C & A & D & A & B & D\\ \hline\hline
    \textbf{Q} & 13 & 14 & 15 & 16 & 17 & 18 & 19 & 20 & 21 & 22 & 23 & 24\\ \hline
    \textbf{Ans} & C & D & B & D & D & C & A & A & C & B & B & D\\ \hline\hline
    \textbf{Q} & 25 & 26 & 27 & 28 & 29 & 30 & 31 & 32 & 33 & 34 & 35 & 36\\ \hline
    \textbf{Ans} & B & C & D & C & A & C & B & A & C & A & D & B\\ \hline\hline
    \textbf{Q} & 37 & 38 & 39 & 40 & 41 & 42 & 43 & 44 & 45 & 46 & 47 & 48\\ \hline
    \textbf{Ans} & A & B & D & D & A & C & D & B & C & B & A & D\\ \hline\hline
    \textbf{Q} & 49 & 50 & 51 & 52 & 53 & 54 & 55 & 56 & 57 & 58 & 59 & 60\\ \hline
    \textbf{Ans} & B & A & A & A & C & C & C & A & C & D & C & A\\ \hline\hline
    \textbf{Q} & 61 & 62 & 63 & 64 & 65 & 66 & 67 & 68 & 69 & 70 & 71 & 72\\ \hline
    \textbf{Ans} & B & B & A & D & B & C & D & B & D & D & B & D\\ \hline
  \end{tabular}
\end{center}
Topic 1 (Projectile Motion): Questions 1--24. \quad Topic 2 (Circular Motion): Questions 25--48.
\quad Topic 3 (Motion in Gravitational Fields): Questions 49--72.

Data used: $g = \qty{9.8}{m.s^{-2}}$, $G = \qty{6.67e-11}{N.m^2.kg^{-2}}$,
$M_\mathrm{E} = \qty{6.0e24}{kg}$, $r_\mathrm{E} = \qty{6.371e6}{m}$.

\begin{practiceanswers}

% =====================================================================
\anstopic{Topic 1 - Projectile Motion}

% ---- Q1
\begin{practiceans}
  \correct{B}
  Gravity acts at every point of the path, including the top. So the acceleration is always
  $g = \qty{9.8}{m.s^{-2}}$ downwards.
  \begin{itemize}
    \item \xmark\ A: gravity pulls down, not up.
    \item \xmark\ C: this is the direction of the \emph{velocity} at the top, not the acceleration.
    \item \xmark\ D: the vertical velocity is zero at the top, but the acceleration is not.
  \end{itemize}
\end{practiceans}

% ---- Q2
\Needspace{9\baselineskip}
\begin{practiceans}
  \correct{A}
  The marble starts with zero vertical velocity, so the time of flight depends only on the height:
  $s_y = \frac{1}{2}gt^2$. A taller bench gives a longer time. The range is $u_x t$. The same $u_x$
  for a longer time gives a bigger range.
  \begin{itemize}
    \item \xmark\ B: the time of flight depends on the height.
    \item \xmark\ C: range $= u_x t$, so a longer time gives a longer range.
    \item \xmark\ D: both change.
  \end{itemize}
\end{practiceans}

% ---- Q3
\Needspace{9\baselineskip}
\begin{practiceans}
  \correct{C}
  At the highest point the vertical velocity is zero. Only the horizontal component is left:
  $v = u\cos\theta = 9.0 \times \cos 35^\circ = \qty{7.4}{m.s^{-1}}$.
  \begin{itemize}
    \item \xmark\ A: only the \emph{vertical} velocity is zero at the top.
    \item \xmark\ B: this is $u\sin\theta$, the vertical component at launch.
    \item \xmark\ D: this is the launch speed. The vertical part is gone at the top.
  \end{itemize}
\end{practiceans}

% ---- Q4
\Needspace{9\baselineskip}
\begin{practiceans}
  \correct{A}
  No horizontal force acts (no air resistance), so the horizontal velocity stays the same.
  \begin{itemize}
    \item \xmark\ B: the vertical component is zero at $Y$, but not at launch.
    \item \xmark\ C and D: the speed and kinetic energy are smaller at $Y$, because the vertical part
          of the velocity is zero there.
  \end{itemize}
\end{practiceans}

% ---- Q5
\Needspace{9\baselineskip}
\begin{practiceans}
  \correct{C}
  The time of flight depends only on the vertical component: $t = \dfrac{2u_y}{g}$. The range is
  $u_x t$. Increasing only $u_x$ keeps $t$ the same and makes the range bigger.
  \begin{itemize}
    \item \xmark\ A: a bigger $u_y$ makes the time of flight longer.
    \item \xmark\ B: a smaller $u_y$ makes the time of flight shorter.
    \item \xmark\ D: a smaller $u_x$ makes the range smaller.
  \end{itemize}
\end{practiceans}

% ---- Q6
\Needspace{9\baselineskip}
\begin{practiceans}
  \correct{B}
  Horizontally the velocity is constant, so the horizontal gaps are equal. Vertically the ball
  starts from rest and speeds up, so the vertical gaps grow (distance fallen $\propto t^2$: the gaps
  are in the ratio 1\,:\,3\,:\,5\,:\,7).
  \begin{itemize}
    \item \xmark\ A: equal vertical gaps mean no vertical acceleration (a straight line).
    \item \xmark\ C: the horizontal gaps get smaller. That needs a horizontal force, such as air
          resistance.
    \item \xmark\ D: the vertical gaps get smaller, but gravity makes the ball fall faster and faster.
  \end{itemize}
\end{practiceans}

% ---- Q7
\Needspace{9\baselineskip}
\begin{practiceans}
  \correct{C}
  The only force is gravity, which is vertical and constant. So the horizontal acceleration is zero
  at every point, and the vertical acceleration is the same (\qty{9.8}{m.s^{-2}} down) at every
  point, on the way up and on the way down.
  \begin{itemize}
    \item \xmark\ A: the horizontal \emph{velocity} is constant, so the horizontal acceleration is zero.
    \item \xmark\ B and D: these vertical arrows show the vertical \emph{velocity} (up, then down). The
          acceleration is downwards even while the ball rises.
  \end{itemize}
\end{practiceans}

% ---- Q8
\Needspace{9\baselineskip}
\begin{practiceans}
  \correct{A}
  No horizontal force acts, so the horizontal component of the velocity does not change: $v_x = u_x$.
  \begin{itemize}
    \item \xmark\ B: the horizontal component does not change.
    \item \xmark\ C: the object is moving downwards, so the vertical component of $v$ is downwards.
    \item \xmark\ D: the object hits the wall \emph{above} its launch height. From $v_y^2 = u_y^2 - 2gh$,
          the vertical speed there is smaller than $u_y$. It equals $u_y$ only back at the launch
          height.
  \end{itemize}
\end{practiceans}

% ---- Q9
\Needspace{9\baselineskip}
\begin{practiceans}
  \correct{D}
  Both balls start with zero vertical velocity, from the same height, with the same $g$. So they
  land at the same time with the same vertical velocity. Y also has a horizontal velocity, so its
  final speed $\sqrt{v_x^2 + v_y^2}$ is greater.
  \begin{itemize}
    \item \xmark\ A: the time to fall depends only on the vertical motion.
    \item \xmark\ B: the horizontal speed does not change the time to fall.
    \item \xmark\ C: the vertical speeds are the same, but Y also has a horizontal speed.
  \end{itemize}
\end{practiceans}

% ---- Q10
\Needspace{9\baselineskip}
\begin{practiceans}
  \correct{A}
  $x = u_x t$, a straight line through the origin. $y = u_y t - \frac{1}{2}gt^2$, a parabola that
  goes up and comes back to zero when the ball lands.
  \begin{itemize}
    \item \xmark\ B: the two graphs are swapped.
    \item \xmark\ C: a straight-line $y$ graph means no gravity.
    \item \xmark\ D: a constant $x$ means no horizontal motion.
  \end{itemize}
\end{practiceans}

% ---- Q11
\Needspace{9\baselineskip}
\begin{practiceans}
  \correct{B}
  With \emph{down} positive, the ball first moves up, so $v$ starts negative. It becomes less
  negative, is zero at the top, then becomes positive. The acceleration is $g$ downwards, which is
  $+\qty{9.8}{m.s^{-2}}$ all the time. So $v$ is a straight line with a positive gradient, and $a$ is
  a constant positive value.
  \begin{itemize}
    \item \xmark\ A: these are the graphs for \emph{up} positive.
    \item \xmark\ C: a rising $v$ graph means a positive acceleration, but this $a$ graph is negative.
    \item \xmark\ D: the V shape is the \emph{speed}, not the velocity.
  \end{itemize}
\end{practiceans}

% ---- Q12
\Needspace{9\baselineskip}
\begin{practiceans}
  \correct{D}
  $E_\mathrm{p} = mgh$. The height rises and falls as a parabola in time. It starts and ends at
  zero, because the ball lands at its launch height.
  \begin{itemize}
    \item \xmark\ A: this is the shape of the kinetic energy graph (largest at the start and end).
    \item \xmark\ B: the height does not keep increasing.
    \item \xmark\ C: $E_\mathrm{p}$ must be zero at the start and end.
  \end{itemize}
\end{practiceans}

% ---- Q13
\Needspace{9\baselineskip}
\begin{practiceans}
  \correct{C}
  From 0 to \qty{2}{s} the velocity is positive (rising). At \qty{2}{s} it is zero (the top:
  $t = \frac{19.6}{9.8} = \qty{2}{s}$).
  From 2 to \qty{4}{s} it is negative (falling). The areas above and below the axis are equal, so
  the projectile is back at the start at \qty{4}{s}. The velocity changes steadily, so the
  displacement--time graph is a parabola, highest at \qty{2}{s}.
  \begin{itemize}
    \item \xmark\ A: straight lines mean constant speed, but the speed changes all the time.
    \item \xmark\ B: this graph is upside down (lowest at \qty{2}{s}).
    \item \xmark\ D: this graph is still rising at \qty{4}{s}, but the velocity is negative after
          \qty{2}{s}.
  \end{itemize}
\end{practiceans}

% ---- Q14
\Needspace{9\baselineskip}
\begin{practiceans}
  \correct{D}
  At the maximum height $v_y = 0$, so $u_y^2 = 2gh = 2 \times 9.8 \times 60$, which gives
  $u_y = \qty{34.3}{m.s^{-1}}$. Then $\sin\theta = \dfrac{34.3}{40}$, so $\theta = 59^\circ$.
  \begin{itemize}
    \item \xmark\ A: uses $\cos\theta = \dfrac{34.3}{40}$.
    \item \xmark\ B: \ang{45} gives the greatest range, not this height.
    \item \xmark\ C: uses $\tan\theta = \dfrac{60}{40}$ (height divided by speed).
  \end{itemize}
\end{practiceans}

% ---- Q15
\Needspace{9\baselineskip}
\begin{practiceans}
  \correct{B}
  The range is $R = \dfrac{u^2\sin 2\theta}{g}$. It is zero at \ang{0} and at \ang{90} (straight up
  and down), and greatest at \ang{45}. The graph is symmetrical about \ang{45}.
  \begin{itemize}
    \item \xmark\ A and C: the range does not keep increasing. At \ang{90} the ball lands where it
          started.
    \item \xmark\ D: the range depends on the angle.
  \end{itemize}
\end{practiceans}

% ---- Q16
\Needspace{9\baselineskip}
\begin{practiceans}
  \correct{D}
  After \qty{1.5}{s}: $v_y = gt = 9.8 \times 1.5 = \qty{14.7}{m.s^{-1}}$ downwards. The velocity is at
  \ang{30} below the horizontal, so $\sin 30^\circ = \dfrac{v_y}{v}$:
  $v = \dfrac{14.7}{0.5} = \qty{29.4}{m.s^{-1}}$. (The horizontal part is
  $\dfrac{14.7}{\tan 30^\circ} = \qty{25.5}{m.s^{-1}}$.)
  \begin{itemize}
    \item \xmark\ A: this is only the vertical component.
    \item \xmark\ B: uses $\dfrac{v_y}{\cos 30^\circ}$, the wrong trigonometric ratio.
    \item \xmark\ C: this is only the horizontal component.
  \end{itemize}
\end{practiceans}

% ---- Q17
\Needspace{9\baselineskip}
\begin{practiceans}
  \correct{D}
  At $Y$, $v_y = 0$ but $v_x$ is not zero. So the kinetic energy is at its smallest, but it is not
  zero.
  \begin{itemize}
    \item \xmark\ A: the horizontal motion continues at $Y$.
    \item \xmark\ B: the gravitational potential energy is a \emph{maximum} at $Y$.
    \item \xmark\ C: with no air resistance, the total mechanical energy stays the same.
  \end{itemize}
\end{practiceans}

% ---- Q18
\Needspace{9\baselineskip}
\begin{practiceans}
  \correct{C}
  After the string is cut, no horizontal force acts on the mass, so it keeps the bus's original
  horizontal velocity. The bus speeds up, so the bus moves forward faster than the mass. Relative
  to the bus, the mass starts from rest and has two constant accelerations: \qty{2}{m.s^{-2}}
  towards the back and \qty{9.8}{m.s^{-2}} down. Starting from rest with a constant acceleration
  gives a straight line, downwards and towards the back.
  \begin{itemize}
    \item \xmark\ A: this would happen only if the bus kept a constant velocity.
    \item \xmark\ B: this would happen if the bus slowed down (as in the 2018 HSC question).
    \item \xmark\ D: a parabola needs a starting velocity at an angle to the acceleration. Here the
          mass starts from rest relative to the bus.
  \end{itemize}
\end{practiceans}

% ---- Q19
\Needspace{9\baselineskip}
\begin{practiceans}
  \correct{A}
  Weight $= mg = 5000 \times 9.8 = \qty{49000}{N}$. The thrust is \qty{40000}{N} upwards, so the net
  force is \qty{9000}{N} downwards. The rocket keeps its horizontal velocity and has a constant
  downward acceleration, so it follows a parabola that curves downwards, like a projectile.
  \begin{itemize}
    \item \xmark\ B: a straight line needs the velocity and the net force in the same line. The
          velocity starts horizontal and the force is vertical.
    \item \xmark\ C and D: the net force is downwards, so the rocket cannot curve upwards.
  \end{itemize}
\end{practiceans}

% ---- Q20
\Needspace{9\baselineskip}
\begin{practiceans}
  \correct{A}
  The lift is moving down and slowing down, so its acceleration is \emph{upwards}. In the lift's
  frame of reference, the projectile then falls with an acceleration bigger than $g$. It falls
  faster, so its path curves down more steeply and it lands closer to the launcher.
  \begin{itemize}
    \item \xmark\ B: a nearly straight path would need a lift accelerating downwards at nearly $g$.
    \item \xmark\ C: a flatter path happens when the lift accelerates \emph{downwards} (for example,
          moving up and slowing down, as in the 2016 HSC question).
    \item \xmark\ D: curving upwards would need the lift to accelerate downwards faster than $g$.
  \end{itemize}
\end{practiceans}

% ---- Q21
\Needspace{9\baselineskip}
\begin{practiceans}
  \correct{C}
  $u_x = 25\cos 40^\circ = \qty{19.2}{m.s^{-1}}$ and $u_y = 25\sin 40^\circ = \qty{16.1}{m.s^{-1}}$.
  Time of flight $t = \dfrac{2u_y}{g} = \qty{3.28}{s}$. Range $= u_x t = 19.2 \times 3.28 = \qty{62.8}{m}$.
  \begin{itemize}
    \item \xmark\ A: uses only the time to reach the top (\qty{1.64}{s}).
    \item \xmark\ B: uses $u\sin\theta$ for the horizontal velocity.
    \item \xmark\ D: this is $\dfrac{u^2}{g}$, the range at \ang{45}, not at \ang{40}.
  \end{itemize}
\end{practiceans}

% ---- Q22
\Needspace{9\baselineskip}
\begin{practiceans}
  \correct{B}
  $u_y = 30\sin 50^\circ = \qty{23.0}{m.s^{-1}}$. At the top $v_y = 0$:
  $h = \dfrac{u_y^2}{2g} = \dfrac{23.0^2}{2 \times 9.8} = \qty{26.9}{m}$.
  \begin{itemize}
    \item \xmark\ A: uses $u\cos\theta$ instead of $u\sin\theta$.
    \item \xmark\ C: uses the whole speed (\qty{30}{m.s^{-1}}) instead of $u_y$.
    \item \xmark\ D: forgets the 2 in $2g$.
  \end{itemize}
\end{practiceans}

% ---- Q23
\Needspace{9\baselineskip}
\begin{practiceans}
  \correct{B}
  Vertically, from rest: $45 = \frac{1}{2} \times 9.8 \times t^2$, so $t = \qty{3.03}{s}$.
  Horizontally: distance $= 12 \times 3.03 = \qty{36.4}{m}$.
  \begin{itemize}
    \item \xmark\ A: uses $h = gt^2$ (forgets the $\frac{1}{2}$).
    \item \xmark\ C: uses $t = 45 \div 12$ (height divided by speed), which has no meaning.
    \item \xmark\ D: rearranges wrongly to $t = \sqrt{2hg}$.
  \end{itemize}
\end{practiceans}

% ---- Q24
\Needspace{9\baselineskip}
\begin{practiceans}
  \correct{D}
  Take up as positive: $u_y = 15\sin 30^\circ = \qty{7.5}{m.s^{-1}}$, $s_y = \qty{-40}{m}$,
  $a = \qty{-9.8}{m.s^{-2}}$. Then $-40 = 7.5t - 4.9t^2$, so $4.9t^2 - 7.5t - 40 = 0$:
  \[ t = \frac{7.5 + \sqrt{7.5^2 + 4(4.9)(40)}}{2(4.9)} = \qty{3.72}{s} \]
  \begin{itemize}
    \item \xmark\ A: this is only the time to come back down to the launch height ($\frac{2u_y}{g}$).
    \item \xmark\ B: treats the vertical velocity as downwards at the start.
    \item \xmark\ C: the time to fall \qty{40}{m} from rest. It ignores the upward start.
  \end{itemize}
\end{practiceans}

% =====================================================================
\anstopic{Topic 2 - Circular Motion}

% ---- Q25
\begin{practiceans}
  \correct{B}
  When the string breaks, no horizontal force acts. By Newton's first law, the puck moves in a
  straight line in the direction of its velocity at $X$. This is along the tangent. At the bottom
  of an anticlockwise circle, the velocity points to the right.
  \begin{itemize}
    \item \xmark\ A: staying on the circle needs a centripetal force.
    \item \xmark\ C: a curved path needs a force.
    \item \xmark\ D: there is no outward (“centrifugal”) force.
  \end{itemize}
\end{practiceans}

% ---- Q26
\Needspace{9\baselineskip}
\begin{practiceans}
  \correct{C}
  $\tau = rF\sin\theta$, so $F = \dfrac{\tau}{r\sin\theta} = \dfrac{90}{0.40 \times \sin 60^\circ}
  = \qty{260}{N}$.
  \begin{itemize}
    \item \xmark\ A: multiplies by $r\sin\theta$ instead of dividing.
    \item \xmark\ B: ignores the angle.
    \item \xmark\ D: uses $\cos 60^\circ$.
  \end{itemize}
\end{practiceans}

% ---- Q27
\Needspace{9\baselineskip}
\begin{practiceans}
  \correct{D}
  $\tau = rF\sin\theta$. Same $r$, same angle and twice the force gives twice the torque.
  \begin{itemize}
    \item \xmark\ A: $\sin 30^\circ = 0.5$, so the torque halves.
    \item \xmark\ B: $\frac{r}{2} \times 2F = rF$, so the torque stays the same.
    \item \xmark\ C: at the hinges $r = 0$, so the torque is zero.
  \end{itemize}
\end{practiceans}

% ---- Q28
\Needspace{9\baselineskip}
\begin{practiceans}
  \correct{C}
  $\tau = rF\sin\theta$. A: $L \times F = FL$. B: $\frac{L}{2} \times 2F = FL$.
  C: $L \times 3F \times \sin 30^\circ = 1.5FL$. D: the force is along the bar ($\theta = 0$), so
  $\tau = 0$.
  \begin{itemize}
    \item \xmark\ A and B: $FL$ is less than $1.5FL$.
    \item \xmark\ D: a force along the bar gives no torque.
  \end{itemize}
\end{practiceans}

% ---- Q29
\Needspace{9\baselineskip}
\begin{practiceans}
  \correct{A}
  Period $T = \dfrac{8.0}{5} = \qty{1.6}{s}$. \ang{150} is $\dfrac{150}{360}$ of a lap, so
  $t = \dfrac{150}{360} \times 1.6 = \qty{0.67}{s}$.
  \begin{itemize}
    \item \xmark\ B: this is for \ang{210}, going the other way around.
    \item \xmark\ C: uses \qty{8.0}{s} as the period.
    \item \xmark\ D: uses $\theta = 150$ (degrees) in $t = \dfrac{\theta}{\omega}$, with $\omega$ in
          \si{rad.s^{-1}}.
  \end{itemize}
\end{practiceans}

% ---- Q30
\Needspace{9\baselineskip}
\begin{practiceans}
  \correct{C}
  $f = \qty{2.0}{Hz}$, so $v = 2\pi rf = 2\pi \times 0.80 \times 2.0 = \qty{10.1}{m.s^{-1}}$.
  $F = \dfrac{mv^2}{r} = \dfrac{0.150 \times 10.1^2}{0.80} = \qty{19}{N}$.
  \begin{itemize}
    \item \xmark\ A: uses $\omega = 2.0$ (forgets $\omega = 2\pi f$).
    \item \xmark\ B: uses $T = \qty{2.0}{s}$ instead of $f = \qty{2.0}{Hz}$.
    \item \xmark\ D: uses the mass in grams.
  \end{itemize}
\end{practiceans}

% ---- Q31
\Needspace{9\baselineskip}
\begin{practiceans}
  \correct{B}
  At the design speed, the horizontal component of the normal force alone gives the centripetal
  force. A faster car needs a bigger centripetal force ($\frac{mv^2}{r}$), and the car tends to slide
  up and out. Friction opposes this, so it acts \emph{down} the slope. Its horizontal component
  points towards the centre and adds to the centripetal force.
  \begin{itemize}
    \item \xmark\ A: friction acts up the slope only when the car is \emph{slower} than the design
          speed (it tends to slide down).
    \item \xmark\ C: the weight $mg$ does not change with speed.
    \item \xmark\ D: there is no outward “centrifugal” force. The net force points towards
          the centre.
  \end{itemize}
\end{practiceans}

% ---- Q32
\Needspace{9\baselineskip}
\begin{practiceans}
  \correct{A}
  $a = \dfrac{v^2}{r}$. Halving $v$ makes $v^2$ 4 times smaller. Tripling $r$ makes $a$ 3 times
  smaller. Together: $\dfrac{a}{4 \times 3} = \dfrac{a}{12}$.
  \begin{itemize}
    \item \xmark\ B: does not square the speed.
    \item \xmark\ C: multiplies by 3 instead of dividing.
    \item \xmark\ D: doubles the speed instead of halving it.
  \end{itemize}
\end{practiceans}

% ---- Q33
\Needspace{9\baselineskip}
\begin{practiceans}
  \correct{C}
  The vehicle moves in a horizontal circle at a constant speed, so the net force is the centripetal
  force: horizontal, towards the centre (to the left).
  \begin{itemize}
    \item \xmark\ A: this is the direction of the normal force, one of the forces, not the net force.
    \item \xmark\ B: a net force down the slope would make the vehicle slide down.
    \item \xmark\ D: this is the direction of the weight.
  \end{itemize}
\end{practiceans}

% ---- Q34
\Needspace{9\baselineskip}
\begin{practiceans}
  \correct{A}
  $\tau = rF\sin\theta$, where $\theta$ is the angle between $r$ and $F$ (\ang{25}):
  $\tau = 0.50 \times 40 \times \sin 25^\circ = \qty{8.5}{N.m}$.
  \begin{itemize}
    \item \xmark\ B: uses $\cos 25^\circ$.
    \item \xmark\ C: ignores the angle.
    \item \xmark\ D: uses $r$ in centimetres.
  \end{itemize}
\end{practiceans}

% ---- Q35
\Needspace{9\baselineskip}
\begin{practiceans}
  \correct{D}
  The hanging mass does not move, so the tension is $Mg$. The tension provides the centripetal
  force on the stopper: $Mg = \dfrac{mv^2}{r}$, so $v = \sqrt{\dfrac{Mgr}{m}}$.
  \begin{itemize}
    \item \xmark\ A: the two masses are swapped.
    \item \xmark\ B: forgets the square root (this is $v^2$).
    \item \xmark\ C: $r$ is on the wrong side of the fraction.
  \end{itemize}
\end{practiceans}

% ---- Q36
\Needspace{9\baselineskip}
\begin{practiceans}
  \correct{B}
  $F = m\omega^2 r$, so $m = \dfrac{F}{\omega^2 r} = \dfrac{3.2}{(3\pi)^2 \times 0.60} = \qty{0.060}{kg}$.
  \begin{itemize}
    \item \xmark\ A: uses $m = \dfrac{Fr}{\omega^2}$ ($r$ on the wrong side).
    \item \xmark\ C: does not square $\omega$.
    \item \xmark\ D: uses $\omega = 3$ (forgets the $\pi$).
  \end{itemize}
\end{practiceans}

% ---- Q37
\Needspace{9\baselineskip}
\begin{practiceans}
  \correct{A}
  In projectile motion the only force is gravity: constant in size and direction. In uniform
  circular motion the net force always points to the centre, so its direction keeps changing.
  \begin{itemize}
    \item \xmark\ B: the acceleration of a projectile is $g$ at every point, including the top.
    \item \xmark\ C: it is the other way around: a projectile's speed changes, and the speed in
          uniform circular motion is constant.
    \item \xmark\ D: also the other way around. Gravity does work on a projectile (its kinetic energy
          changes). In uniform circular motion the net force is at \ang{90} to the velocity, so it
          does no work.
  \end{itemize}
\end{practiceans}

% ---- Q38
\Needspace{9\baselineskip}
\begin{practiceans}
  \correct{B}
  $v = \dfrac{72}{3.6} = \qty{20}{m.s^{-1}}$. $\tan\theta = \dfrac{v^2}{rg} = \dfrac{20^2}{120 \times
  9.8} = 0.340$, so $\theta = 19^\circ$.
  \begin{itemize}
    \item \xmark\ A: does not square $v$.
    \item \xmark\ C: the fraction is upside down.
    \item \xmark\ D: uses \qty{72}{km.h^{-1}} without changing it to \si{m.s^{-1}}.
  \end{itemize}
\end{practiceans}

% ---- Q39
\Needspace{9\baselineskip}
\begin{practiceans}
  \correct{D}
  $F = \dfrac{mv^2}{r}$ and $v = \dfrac{2\pi r}{T}$, so $F = \dfrac{4\pi^2 mr}{T^2}$. With $m$ and
  $r$ fixed, $F \propto \dfrac{1}{T^2}$. So $F$ against $\dfrac{1}{T^2}$ is a straight line through
  the origin.
  \begin{itemize}
    \item \xmark\ A, B and C: these pairs give curves.
  \end{itemize}
\end{practiceans}

% ---- Q40
\Needspace{9\baselineskip}
\begin{practiceans}
  \correct{D}
  On the straight parts there is no sideways acceleration (zero). In the left turn, the centripetal
  acceleration points to the centre of the turn, which is on the car's left: a constant positive
  value. In the right turn, the centre is on the car's right: a constant negative value.
  \begin{itemize}
    \item \xmark\ A: the signs are swapped.
    \item \xmark\ B: the right turn must give a negative reading.
    \item \xmark\ C: the acceleration lasts for the whole turn, not just a moment at each change.
  \end{itemize}
\end{practiceans}

% ---- Q41
\Needspace{9\baselineskip}
\begin{practiceans}
  \correct{A}
  Vertically: $T\cos\theta = mg$. Horizontally: $T\sin\theta = \dfrac{mv^2}{r}$. Dividing:
  $\tan\theta = \dfrac{v^2}{rg}$, so $v = \sqrt{rg\tan\theta}$.
  \begin{itemize}
    \item \xmark\ B: $\tan\theta$ is upside down.
    \item \xmark\ C: uses $\sin\theta$ instead of $\tan\theta$.
    \item \xmark\ D: forgets the square root.
  \end{itemize}
\end{practiceans}

% ---- Q42
\Needspace{9\baselineskip}
\begin{practiceans}
  \correct{C}
  The net force towards the centre must be $\dfrac{mv^2}{r}$ (constant). At the top, the weight
  points towards the centre and helps: $T + mg = \dfrac{mv^2}{r}$, so $T$ is smallest. At the
  bottom, the weight points away from the centre: $T - mg = \dfrac{mv^2}{r}$, so $T$ is largest.
  $T$ rises from the top to a maximum at the bottom (half a turn), then falls again.
  \begin{itemize}
    \item \xmark\ A: a constant tension is true for a horizontal circle, not a vertical one.
    \item \xmark\ B: this is upside down: the tension is smallest at the top.
    \item \xmark\ D: at the sides the weight is along the path, so $T = \dfrac{mv^2}{r}$ there, a
          middle value, not a maximum.
  \end{itemize}
\end{practiceans}

% ---- Q43
\Needspace{9\baselineskip}
\begin{practiceans}
  \correct{D}
  The puck moves in a circle, so the net force points to the centre. The weight is balanced by the
  table. The only horizontal force is the normal force from the wall, so it provides the
  centripetal force.
  \begin{itemize}
    \item \xmark\ A: the direction of the velocity changes, so the puck accelerates. The net force
          is not zero.
    \item \xmark\ B: no force is needed along the motion, because the speed is constant.
    \item \xmark\ C: there is no “centrifugal” force acting on the puck.
  \end{itemize}
\end{practiceans}

% ---- Q44
\Needspace{9\baselineskip}
\begin{practiceans}
  \correct{B}
  At the minimum speed, the track pushes with zero force at the top, so the weight alone gives the
  centripetal force: $mg = \dfrac{mv^2}{r}$, so $v = \sqrt{gr} = \sqrt{9.8 \times 8.0} =
  \qty{8.9}{m.s^{-1}}$.
  \begin{itemize}
    \item \xmark\ A: uses $\sqrt{\dfrac{g}{r}}$.
    \item \xmark\ C: uses $\sqrt{2gr}$.
    \item \xmark\ D: forgets the square root.
  \end{itemize}
\end{practiceans}

% ---- Q45
\Needspace{9\baselineskip}
\begin{practiceans}
  \correct{C}
  Friction provides the centripetal force: $F = \dfrac{mv^2}{r}$, so
  $v = \sqrt{\dfrac{Fr}{m}} = \sqrt{\dfrac{9600 \times 50}{1200}} = \qty{20}{m.s^{-1}}$.
  \begin{itemize}
    \item \xmark\ A: uses $\sqrt{\dfrac{F}{mr}}$.
    \item \xmark\ B: this is $\dfrac{F}{m}$, the acceleration in \si{m.s^{-2}}.
    \item \xmark\ D: forgets the square root (this is $v^2$).
  \end{itemize}
\end{practiceans}

% ---- Q46
\Needspace{9\baselineskip}
\begin{practiceans}
  \correct{B}
  $f = \dfrac{120}{60} = \qty{2.0}{rev.s^{-1}}$, so $\omega = 2\pi f = \qty{12.6}{rad.s^{-1}}$.
  $v = \omega r = 12.6 \times 0.30 = \qty{3.8}{m.s^{-1}}$.
  \begin{itemize}
    \item \xmark\ A: uses $\omega = 2.0$ (forgets the $2\pi$).
    \item \xmark\ C: uses 120 as $\omega$.
    \item \xmark\ D: uses $f = 120$, which is per minute, not per second.
  \end{itemize}
\end{practiceans}

% ---- Q47
\Needspace{9\baselineskip}
\begin{practiceans}
  \correct{A}
  The tension is always at \ang{90} to the velocity, so it does no work ($W = Fs\cos 90^\circ = 0$).
  This also matches the speed staying constant: the kinetic energy does not change.
  \begin{itemize}
    \item \xmark\ B: this is the kinetic energy of the mass, not work done.
    \item \xmark\ C: this is $F \times r$.
    \item \xmark\ D: this is $F \times$ circumference, but the force is at \ang{90} to the motion.
  \end{itemize}
\end{practiceans}

% ---- Q48
\Needspace{9\baselineskip}
\begin{practiceans}
  \correct{D}
  $\tau = rF\sin\theta$, where $\theta$ is the angle between the crank and the (vertical) force.
  When the crank is horizontal, $\theta = 90^\circ$ and the torque is greatest.
  \begin{itemize}
    \item \xmark\ A and C: the force is along the crank ($\theta = 0$ or \ang{180}), so the torque is
          zero.
    \item \xmark\ B: $\sin 45^\circ = 0.71$, so the torque is only 71\% of the maximum.
  \end{itemize}
\end{practiceans}

% =====================================================================
\anstopic{Topic 3 - Motion in Gravitational Fields}

% ---- Q49
\begin{practiceans}
  \correct{B}
  Newton's third law: the two forces are equal in size and opposite in direction. ($F =
  \dfrac{Gm(4m)}{d^2}$ for both.)
  \begin{itemize}
    \item \xmark\ A, C and D: the two forces of a third-law pair are always equal in size.
  \end{itemize}
\end{practiceans}

% ---- Q50
\Needspace{9\baselineskip}
\begin{practiceans}
  \correct{A}
  $g = \dfrac{GM}{r^2}$, so $g$ changes by $\dfrac{5}{2.5^2} = \dfrac{5}{6.25} = 0.8$.
  \begin{itemize}
    \item \xmark\ B: divides by $r$, not $r^2$.
    \item \xmark\ C: ignores the change in radius.
    \item \xmark\ D: multiplies by $r^2$ instead of dividing.
  \end{itemize}
\end{practiceans}

% ---- Q51
\Needspace{9\baselineskip}
\begin{practiceans}
  \correct{A}
  $F = \dfrac{GMm}{r^2}$: one mass $\times 2$ and $r \times 3$ gives $\dfrac{2}{3^2} = \dfrac{2}{9}$.
  \begin{itemize}
    \item \xmark\ B: divides by $r$, not $r^2$.
    \item \xmark\ C: multiplies by the change in $r$ instead of dividing.
    \item \xmark\ D: multiplies by $r^2$ instead of dividing.
  \end{itemize}
\end{practiceans}

% ---- Q52
\Needspace{9\baselineskip}
\begin{practiceans}
  \correct{A}
  $r = \num{6.371e6} + \num{1.0e6} = \qty{7.371e6}{m}$.
  $g = \dfrac{GM}{r^2} = \dfrac{(\num{6.67e-11})(\num{6.0e24})}{(\num{7.371e6})^2} = \qty{7.4}{N.kg^{-1}}$.
  \begin{itemize}
    \item \xmark\ B: uses $g \propto \dfrac{1}{r}$ (no square).
    \item \xmark\ C: uses only the altitude as $r$.
    \item \xmark\ D: forgets to square $r$.
  \end{itemize}
\end{practiceans}

% ---- Q53
\Needspace{9\baselineskip}
\begin{practiceans}
  \correct{C}
  $v = \sqrt{\dfrac{GM}{r}}$, so $v$ gets smaller as $r$ gets bigger. It is a curve that never
  reaches zero.
  \begin{itemize}
    \item \xmark\ A: higher orbits are slower, not faster.
    \item \xmark\ B: the speed depends on $r$.
    \item \xmark\ D: the shape is a curve ($v \propto \frac{1}{\sqrt{r}}$), not a straight line to zero.
  \end{itemize}
\end{practiceans}

% ---- Q54
\Needspace{9\baselineskip}
\begin{practiceans}
  \correct{C}
  In astronomical units and years, $r^3 = T^2$ for everything orbiting the Sun. So
  $r = 1.88^{2/3} = 1.52$.
  \begin{itemize}
    \item \xmark\ A: the cube root of $T$ ($r^3 = T$).
    \item \xmark\ B: the square root of $T$ ($r^2 = T$).
    \item \xmark\ D: $r$ and $T$ swapped ($r^2 = T^3$).
  \end{itemize}
\end{practiceans}

% ---- Q55
\Needspace{9\baselineskip}
\begin{practiceans}
  \correct{C}
  $r = \num{6.371e6} + \num{8.0e5} = \qty{7.171e6}{m}$.
  $T = 2\pi\sqrt{\dfrac{r^3}{GM}} = \qty{6.03e3}{s} = 101$ minutes.
  \begin{itemize}
    \item \xmark\ A: uses only the altitude as $r$.
    \item \xmark\ B: uses $r = r_\mathrm{E}$ and forgets the altitude.
    \item \xmark\ D: forgets the square root.
  \end{itemize}
\end{practiceans}

% ---- Q56
\Needspace{9\baselineskip}
\begin{practiceans}
  \correct{A}
  A geostationary satellite has the same period as Earth's rotation (24 hours), so it stays above
  one point: it can watch one region all the time. A low Earth orbit satellite is much closer to
  the surface, so its photographs show more detail.
  \begin{itemize}
    \item \xmark\ B: geostationary satellites are about \qty{36000}{km} high, much further away.
    \item \xmark\ C: low Earth orbit satellites take about 90 minutes per orbit, so they move quickly
          across the sky.
    \item \xmark\ D: their periods are very different (about 90 minutes and 24 hours).
  \end{itemize}
\end{practiceans}

% ---- Q57
\Needspace{9\baselineskip}
\begin{practiceans}
  \correct{C}
  Kepler's third law: $T^2 \propto r^3$, with the same constant for all the planets of the Sun.
  \begin{itemize}
    \item \xmark\ A: the first law says the orbits are ellipses, with the Sun at one focus.
    \item \xmark\ B: the second law means the speed changes (fastest near the Sun).
    \item \xmark\ D: $\dfrac{r^3}{T^2} = \dfrac{GM_\mathrm{Sun}}{4\pi^2}$ does not contain the planet's
          mass.
  \end{itemize}
\end{practiceans}

% ---- Q58
\Needspace{9\baselineskip}
\begin{practiceans}
  \correct{D}
  The gravitational force always points towards the centre of Earth, at every point of every
  orbit, including elliptical orbits.
  \begin{itemize}
    \item \xmark\ A: gravity pulls towards Earth, not along the path.
    \item \xmark\ B: the centre of the ellipse is empty space. Nothing there pulls on the satellite.
    \item \xmark\ C: the force is at \ang{90} to the path only in a circular orbit.
  \end{itemize}
\end{practiceans}

% ---- Q59
\Needspace{9\baselineskip}
\begin{practiceans}
  \correct{C}
  \[ M = \frac{4\pi^2 r^3}{GT^2} = \frac{4\pi^2(\num{1.22e9})^3}{(\num{6.67e-11})(\num{1.38e6})^2}
       = \qty{5.6e26}{kg} \]
  \begin{itemize}
    \item \xmark\ A: uses $r^2$ instead of $r^3$.
    \item \xmark\ B: forgets the $4\pi^2$.
    \item \xmark\ D: forgets to square $T$.
  \end{itemize}
\end{practiceans}

% ---- Q60
\Needspace{9\baselineskip}
\begin{practiceans}
  \correct{A}
  An altitude of $3r_\mathrm{E}$ means $r = 4r_\mathrm{E}$ from the centre. $F \propto \dfrac{1}{r^2}$,
  so $F = \dfrac{800}{4^2} = \qty{50}{N}$.
  \begin{itemize}
    \item \xmark\ B: uses $r = 3r_\mathrm{E}$ (forgets to add Earth's radius).
    \item \xmark\ C: uses $F \propto \dfrac{1}{r}$.
    \item \xmark\ D: makes both mistakes.
  \end{itemize}
\end{practiceans}

% ---- Q61
\Needspace{9\baselineskip}
\begin{practiceans}
  \correct{B}
  $r = \num{6.371e6} + \num{2.02e7} = \qty{2.657e7}{m}$.
  \[ U = -\frac{GMm}{r} = -\frac{(\num{6.67e-11})(\num{6.0e24})(500)}{\num{2.657e7}} = \qty{-7.5e9}{J} \]
  \begin{itemize}
    \item \xmark\ A: this is the total energy, $\frac{1}{2}U$.
    \item \xmark\ C: uses only the altitude as $r$.
    \item \xmark\ D: uses $r = r_\mathrm{E}$ and forgets the altitude.
  \end{itemize}
\end{practiceans}

% ---- Q62
\Needspace{9\baselineskip}
\begin{practiceans}
  \correct{B}
  In a circular orbit $\dfrac{mv^2}{r} = \dfrac{GMm}{r^2}$, so
  \[ K = \tfrac{1}{2}mv^2 = \frac{GMm}{2r} = \frac{(\num{6.67e-11})(\num{6.0e24})(2000)}{2 \times \num{8.0e6}}
     = \qty{5.0e10}{J} \]
  \begin{itemize}
    \item \xmark\ A: this is $\dfrac{GMm}{r^2}$, the force in newtons, not an energy.
    \item \xmark\ C: this is $\dfrac{GMm}{r}$, the size of $U$ (it forgets the $\frac{1}{2}$).
    \item \xmark\ D: kinetic energy is never negative. This is the total energy.
  \end{itemize}
\end{practiceans}

% ---- Q63
\Needspace{9\baselineskip}
\begin{practiceans}
  \correct{A}
  $K = \dfrac{GMm}{2r}$: when $r$ doubles, $K$ halves (decreases). $E = -\dfrac{GMm}{2r}$: when $r$
  doubles, $E$ becomes less negative (increases). The rocket engine must supply this energy.
  \begin{itemize}
    \item \xmark\ B: a higher orbit is slower, so $K$ decreases, and $E$ increases.
    \item \xmark\ C: $E$ becomes less negative, which is an increase.
    \item \xmark\ D: $K$ decreases.
  \end{itemize}
\end{practiceans}

% ---- Q64
\Needspace{9\baselineskip}
\begin{practiceans}
  \correct{D}
  $U = -\dfrac{GMm}{r}$: it is always negative. It gets closer to zero as $r$ increases, and more
  negative closer to the planet.
  \begin{itemize}
    \item \xmark\ A and C: $U$ is negative everywhere (zero only at infinity).
    \item \xmark\ B: $U \propto -\frac{1}{r}$, not $-r$. This line gets more negative further away.
  \end{itemize}
\end{practiceans}

% ---- Q65
\Needspace{9\baselineskip}
\begin{practiceans}
  \correct{B}
  $v_\mathrm{esc} = \sqrt{\dfrac{2GM}{r}} = \sqrt{\dfrac{2(\num{6.67e-11})(\num{7.35e22})}{\num{1.74e6}}}
  = \qty{2.37e3}{m.s^{-1}}$.
  \begin{itemize}
    \item \xmark\ A: this is the orbital speed $\sqrt{\dfrac{GM}{r}}$ (it forgets the 2).
    \item \xmark\ C: uses $r$ in kilometres.
    \item \xmark\ D: forgets the square root.
  \end{itemize}
\end{practiceans}

% ---- Q66
\Needspace{9\baselineskip}
\begin{practiceans}
  \correct{C}
  $v_\mathrm{esc} = \sqrt{\dfrac{2GM}{r}}$. Same $r$ and 4 times the mass: $v_\mathrm{esc}$ is
  multiplied by $\sqrt{4} = 2$, giving \qty{22.4}{km.s^{-1}}.
  \begin{itemize}
    \item \xmark\ A: divides instead of multiplying.
    \item \xmark\ B: the escape velocity depends on the planet's mass.
    \item \xmark\ D: forgets the square root.
  \end{itemize}
\end{practiceans}

% ---- Q67
\Needspace{9\baselineskip}
\begin{practiceans}
  \correct{D}
  $v_\mathrm{esc} = \sqrt{\dfrac{2GM}{r}}$, so a bigger $r$ gives a smaller escape velocity.
  \begin{itemize}
    \item \xmark\ A: the rocket's mass cancels out, so it does not matter.
    \item \xmark\ B: escape velocity comes from energy, so it works in any direction (as long as the
          object does not hit the planet).
    \item \xmark\ C: escape velocity is $\sqrt{2}$ times the orbital speed at the same distance.
  \end{itemize}
\end{practiceans}

% ---- Q68
\Needspace{9\baselineskip}
\begin{practiceans}
  \correct{B}
  $r_1 = \num{6.371e6} + \num{4.0e5} = \qty{6.771e6}{m}$ and $r_2 = \qty{4.23e7}{m}$.
  \[ \Delta E = \frac{GMm}{2}\left(\frac{1}{r_1} - \frac{1}{r_2}\right) \]
  \[ \Delta E = \frac{(\num{6.67e-11})(\num{6.0e24})(1000)}{2}\left(\frac{1}{\num{6.771e6}} - \frac{1}{\num{4.23e7}}\right)
     = \qty{2.5e10}{J} \]
  \begin{itemize}
    \item \xmark\ A: this is only the size of the total energy in the new orbit.
    \item \xmark\ C: this is only the size of the total energy in the first orbit.
    \item \xmark\ D: this is the change in $U$ only. The satellite also loses kinetic energy.
  \end{itemize}
\end{practiceans}

% ---- Q69
\Needspace{9\baselineskip}
\begin{practiceans}
  \correct{D}
  At the point, $\dfrac{GM_\mathrm{E}m}{d_\mathrm{E}^2} = \dfrac{GM_\mathrm{M}m}{d_\mathrm{M}^2}$, so
  $\dfrac{d_\mathrm{E}}{d_\mathrm{M}} = \sqrt{\dfrac{M_\mathrm{E}}{M_\mathrm{M}}}$. A heavier Moon pulls
  harder, so the balance point moves away from the Moon, closer to Earth.
  \begin{itemize}
    \item \xmark\ A: this is the wrong way.
    \item \xmark\ B: the position depends on the ratio of the masses.
    \item \xmark\ C: between the two bodies the forces always point in opposite directions, so a
          balance point always exists.
  \end{itemize}
\end{practiceans}

% ---- Q70
\Needspace{9\baselineskip}
\begin{practiceans}
  \correct{D}
  The centripetal acceleration of a satellite is $\dfrac{GM}{r^2}$, so $a \propto \dfrac{1}{r^2}$.
  With $r_\mathrm{Q} = 4r_\mathrm{P}$, the ratio is $4^2 = 16$.
  \begin{itemize}
    \item \xmark\ A: this is the ratio of the orbital speeds ($v \propto \frac{1}{\sqrt{r}}$).
    \item \xmark\ B: uses $a \propto \frac{1}{r}$.
    \item \xmark\ C: this is the ratio of the periods ($T \propto r^{3/2}$).
  \end{itemize}
\end{practiceans}

% ---- Q71
\Needspace{9\baselineskip}
\begin{practiceans}
  \correct{B}
  $r = \sqrt[3]{\dfrac{GMT^2}{4\pi^2}} = \qty{4.23e7}{m}$ with $T = \qty{86400}{s}$. Then
  $v = \dfrac{2\pi r}{T} = \dfrac{2\pi(\num{4.23e7})}{86\,400} = \qty{3.1e3}{m.s^{-1}}$.
  \begin{itemize}
    \item \xmark\ A: this is the speed of a point on the equator, on Earth's surface.
    \item \xmark\ C: this is the speed of a satellite in low Earth orbit.
    \item \xmark\ D: this is the escape velocity from Earth's surface.
  \end{itemize}
\end{practiceans}

% ---- Q72
\Needspace{9\baselineskip}
\begin{practiceans}
  \correct{D}
  In a circular orbit $K = \dfrac{GMm}{2r}$ and $U = -\dfrac{GMm}{r}$, so $U = -2K = \qty{-6.0e9}{J}$.
  Then $E = K + U = 3.0 - 6.0 = \qty{-3.0e9}{J}$ (which is $-K$).
  \begin{itemize}
    \item \xmark\ A: $U$ is not $-K$.
    \item \xmark\ B: $E = K + U = \qty{-3.0e9}{J}$, not $+3.0$.
    \item \xmark\ C: $U$ is not $-\frac{1}{2}K$, and the total energy of an orbit is negative.
  \end{itemize}
\end{practiceans}

\end{practiceanswers}
\end{document}
